\section{Invertible sheaves and morphisms into relative Proj}
\label{section-invertible-relative-proj}

\noindent
It seems that we may need the following lemma somewhere.
The situation is the following:
\begin{enumerate}
\item Let $S$ be a scheme.
\item Let $\mathcal{A}$ be a quasi-coherent graded $\mathcal{O}_S$-algebra.
\item Denote $\pi : \underline{\text{Proj}}_S(\mathcal{A}) \to S$ the relative
homogeneous spectrum over $S$.
\item Let $f : X \to S$ be a morphism of schemes.
\item Let $\mathcal{L}$ be an invertible $\mathcal{O}_X$-module.
\item Let $\psi : f^*\mathcal{A} \to
\bigoplus_{d \geq 0} \mathcal{L}^{\otimes d}$
be a homomorphism of graded $\mathcal{O}_X$-algebras.
\end{enumerate}
Given this data set
$$
U(\psi) = \bigcup\nolimits_{(U, V, a)} U_{\psi(a)}
$$
where $(U, V, a)$ satisfies:
\begin{enumerate}
\item $V \subset S$ affine open,
\item $U = f^{-1}(V)$, and
\item $a \in \mathcal{A}(V)_{+}$ is homogeneous.
\end{enumerate}
Namely, then $\psi(a) \in \Gamma(U, \mathcal{L}^{\otimes \deg(a)})$
and $U_{\psi(a)}$ is the corresponding open (see
Modules, Lemma \ref{modules-lemma-s-open}).

\begin{lemma}
\label{lemma-invertible-map-into-relative-proj}
With assumptions and notation as above. The morphism
$\psi$ induces a canonical morphism of schemes over $S$
$$
r_{\mathcal{L}, \psi} :
U(\psi) \longrightarrow \underline{\text{Proj}}_S(\mathcal{A})
$$
together with a map of graded $\mathcal{O}_{U(\psi)}$-algebras
$$
\theta :
r_{\mathcal{L}, \psi}^*\left(
\bigoplus\nolimits_{d \geq 0}
\mathcal{O}_{\underline{\text{Proj}}_S(\mathcal{A})}(d)
\right)
\longrightarrow
\bigoplus\nolimits_{d \geq 0} \mathcal{L}^{\otimes d}|_{U(\psi)}
$$
characterized by the following properties:
\begin{enumerate}
\item For every open $V \subset S$ and every $d \geq 0$ the diagram
$$
\xymatrix{
\mathcal{A}_d(V) \ar[d]_{\psi} \ar[r]_{\psi} &
\Gamma(f^{-1}(V), \mathcal{L}^{\otimes d}) \ar[d]^{restrict} \\
\Gamma(\pi^{-1}(V),
\mathcal{O}_{\underline{\text{Proj}}_S(\mathcal{A})}(d)) \ar[r]^{\theta} &
\Gamma(f^{-1}(V) \cap U(\psi), \mathcal{L}^{\otimes d})
}
$$
is commutative.
\item For any $d \geq 1$ and any open subscheme $W \subset X$
such that $\psi|_W : f^*\mathcal{A}_d|_W \to \mathcal{L}^{\otimes d}|_W$
is surjective the restriction of the morphism $r_{\mathcal{L}, \psi}$
agrees with the morphism $W \to \underline{\text{Proj}}_S(\mathcal{A})$
which exists by the construction of the relative homogeneous spectrum,
see Definition \ref{definition-relative-proj}.
\item For any affine open $V \subset S$, the restriction
$$
(U(\psi) \cap f^{-1}(V), r_{\mathcal{L}, \psi}|_{U(\psi) \cap f^{-1}(V)},
\theta|_{U(\psi) \cap f^{-1}(V)})
$$
agrees via $i_V$ (see Lemma \ref{lemma-glue-relative-proj}) with the triple
$(U(\psi'), r_{\mathcal{L}, \psi'}, \theta')$
of Lemma \ref{lemma-invertible-map-into-proj} associated to the map
$\psi' : A = \mathcal{A}(V) \to \Gamma_*(f^{-1}(V), \mathcal{L}|_{f^{-1}(V)})$
induced by $\psi$.
\end{enumerate}
\end{lemma}

\begin{proof}
Use characterization (3) to construct the morphism $r_{\mathcal{L}, \psi}$
and $\theta$ locally over $S$. Use the uniqueness of
Lemma \ref{lemma-invertible-map-into-proj}
to show that the construction glues. Details omitted.
\end{proof}







